\chapter*{缩略语与符号}
\addcontentsline{toc}{chapter}{缩略语与符号}
\markboth{缩略语与符号}{缩略语与符号}

本节列出常用的缩略语与符号。
\begin{longtable}{@{}p{.23\textwidth}p{.73\textwidth}@{}}
$A$ & 处理或暴露。 \\
A-IPCW & 增广逆删失概率加权。 \\
A-IPW & 增广逆概率加权。 \\
$C$ & 删失。 \\
CRT & 社区随机试验。 \\
i.i.d. & 独立同分布。 \\
IPCW & 逆删失概率加权。 \\
IPW & 逆概率加权。 \\
LTMLE & 纵向靶向极大似然估计／估计量。 \\
MLE & g 公式的极大似然代入估计量。\emph{不要与采用极大似然估计的非代入估计量混淆。MLE 也称为 g 计算。} \\
MSE & 均方误差。 \\
$O$ & 按时间排序的观测数据结构。 \\
$P$ & 可能的数据生成分布。 \\
$p$ & 数据生成分布 $P_0$ 的可能密度。 \\
$P_0$ & 真实数据生成分布；$O\sim P_0$。 \\
$p_0$ & 数据生成分布 $P_0$ 的真实密度。 \\
$P_n$ & 经验概率分布，在每个观测 $O_i$（$i=1,\ldots,n$）上赋予概率 $1/n$。 \\
RCT & 随机对照试验。 \\
SCM & 结构因果模型。 \\
SE & 标准误。 \\
SL & 超级学习器。 \\
TMLE & 靶向极大似然估计／估计量。 \\
$W$ & 协变量向量。 \\
$Y$ & 结局。 \\
$Y_1,Y_0$ & 处理 $A$ 为二元变量时的反事实结局。 \\
\end{longtable}

大写字母表示随机变量，小写字母表示该变量的某个具体取值。如果 $O$ 离散，则 $p_0(o)=P_0(O=o)$ 表示 $O$ 等于 $o$ 的概率；如果 $O$ 连续，则 $p_0(o)$ 表示 $P_0$ 在 $o$ 处的 Lebesgue 密度。为了简化表述，我们经常把 $O$ 当作离散变量，从而将 $P_0(O=o)$ 称为概率。举一个简单例子，假设数据结构为 $O=(W,A,Y)\sim P_0$，且 $O$ 离散。那么，对每个可能的取值 $(w,a,y)$，$p_0(w,a,y)$ 表示 $(W,A,Y)$ 等于 $(w,a,y)$ 的概率。
\begin{longtable}{@{}p{.23\textwidth}p{.73\textwidth}@{}}
$\mathcal M$ & 统计模型，即 $P_0$ 的所有可能概率分布构成的集合。 \\
$P_0\in\mathcal M$ & 已知 $P_0$ 属于统计模型 $\mathcal M$。 \\
\end{longtable}

本书常用“半参数”一词同时涵盖非参数和半参数情形。当半参数不包括非参数、并且需要作额外假设时，我们会明确说明。一个\emph{统计模型}可以加入额外的非统计假设，例如因果假设，以获得更丰富的解释；对于某种参数化 $\theta\mapsto P_\theta$，这类模型通常写为 $\{P_\theta:\theta\in\Theta\}$。我们将其称为\emph{模型}。例如，观测数据 $O=(W,A,Y=Y_A)$ 的概率分布，可以表示为反事实结局 $Y_0,Y_1$ 上、以 $A$ 为缺失指示变量的缺失数据结构。这样，$O$ 的概率分布就由 $(W,Y_0,Y_1)$ 的概率分布，以及给定 $(W,Y_0,Y_1)$ 时处理 $A$ 的条件分布来索引。
\begin{longtable}{@{}p{.23\textwidth}p{.73\textwidth}@{}}
$X=(X_j:j)$ & 内生变量集合，$j=1,\ldots,J$。 \\
$U=(U_{X_j}:j)$ & 外生变量集合。 \\
$P_{U,X}$ & $(U,X)$ 的概率分布。 \\
$p_{U,X}$ & $(U,X)$ 的密度。 \\
$\mathrm{Pa}(X_j)$ & $X$ 中 $X_j$ 的父节点。 \\
$f_{X_j}$ & 关于 $\mathrm{Pa}(X_j)$ 与 $X_j$ 所对应的一个内生变量 $U_{X_j}$ 的函数。 \\
$f=(f_{X_j}:j)$ & 定义结构因果模型的函数 $f_{X_j}$ 的集合。 \\
$\mathcal M^F$ & 结构因果模型所描述的所有可能 $P_{U,X}$ 的集合；包含依据真实知识提出的不可检验假设，即在 $\mathcal M$ 中加入已知成立的额外非统计假设。 \\
$\mathcal M^{F*}$ & 为使目标参数可识别，可能需要进一步加入因果假设后的模型。 \\
$P\to\Psi(P)$ & 目标参数，作为把 $P$ 映射到其相应参数值的映射。 \\
$\Psi(P_0)$ & 真实目标参数。 \\
$\widehat\Psi(P_n)$ & 估计量，作为把经验分布 $P_n$ 映射到其相应值的映射。 \\
$\psi_0=\Psi(P_0)$ & 真实目标参数值。 \\
$\psi_n$ & $\psi_0$ 的估计。 \\
\end{longtable}
考虑 $O=(L_0,A_0,\ldots,L_K,A_K,L_{K+1})\sim P_0$。
\begin{longtable}{@{}p{.23\textwidth}p{.73\textwidth}@{}}
$L_k$ & $t=k$ 时可能随时间变化的协变量；也记作 $L(k)$。 \\
$A_k$ & $t=k$ 时的时变干预节点，可以同时包含处理和删失。 \\
$\mathrm{Pa}(L_k)$ & $=(\bar A_{k-1},\bar L_{k-1})$。 \\
$\mathrm{Pa}(A_k)$ & $=(\bar A_{k-1},\bar L_k)$。 \\
$P_{0,L_k}$ & $P_0$ 下给定 $\mathrm{Pa}(L_k)$ 时 $L_k$ 的真实条件概率分布。 \\
$P_{L_k}$ & $P$ 下给定 $\mathrm{Pa}(L_k)$ 时 $L_k$ 的条件概率分布。 \\
$P_{n,L_k}$ & $L_k$ 的条件概率分布 $P_{0,L_k}$ 的估计。 \\
$P_{0,A_k}$ & $P_0$ 下给定 $\mathrm{Pa}(A_k)$ 时 $A_k$ 的真实条件概率分布。 \\
$P_{A_k}$ & $P$ 下给定 $\mathrm{Pa}(A_k)$ 时 $A_k$ 的条件概率分布。 \\
$P_{n,A_k}$ & $P_0$ 的估计量 $P_n$ 下，给定 $\mathrm{Pa}(A_k)$ 时 $A_k$ 的条件概率分布。 \\
$\epsilon$ & 扰动参数。 \\
$\epsilon_n$ & $\epsilon$ 的估计。 \\
$\{P_\epsilon:\epsilon\}\subset\mathcal M$ & 穿过 $P$ 的子模型。 \\
$H^*$ & 巧妙协变量。 \\
$H_n^*$ & $H^*$ 的估计。 \\
$D(\psi)(O)$ & 关于数据结构 $O$ 和参数的估计函数；简记为 $D(\psi)$。 \\
$D_0^*(O)$ & 有效影响曲线，即典范梯度；也记作 $D^*(P_0)(O)$、$D^*(P_0)$ 或 $D^*(O)$。 \\
$\mathrm{IC}_0(O)$ & 某个估计量在 $P_0$ 处的影响曲线，是 $O$ 的函数。 \\
$\mathrm{IC}_n(O)$ & 影响曲线的估计。 \\
\end{longtable}

在许多章节中，我们重点研究前面介绍的一般数据结构 $O=(L_0,A_0,\ldots,L_K,A_K,L_{K+1})\sim P_0$。在这一情形下，采用下列特定符号定义：
\begin{longtable}{@{}p{.23\textwidth}p{.73\textwidth}@{}}
$L_0$ & 基线协变量。 \\
$\bar L$ & $=(L_0,\ldots,L_{K+1})$。 \\
$\bar A$ & $=(A_0,\ldots,A_K)$。 \\
$L_d$ & 策略 $d$ 下的反事实结局。 \\
$d_0$ & 依赖于 $P_0$ 的最优策略。 \\
$Q_{0,L_k}$ & $L_k$ 的真实条件概率分布。 \\
$Q_{L_k}$ & $L_k$ 的可能条件概率分布。 \\
$Q_{n,L_k}$ & $Q_{0,L_k}$ 的估计。 \\
$Q$ & $=(Q_{L_0},\ldots,Q_{L_{K+1}})$。 \\
$L(O,Q)$ & 一个损失函数的例子，是 $O$ 与 $Q$ 的函数；也记作 $L(Q)(O)$ 或 $L(Q)$。 \\
$\{Q_\epsilon:\epsilon\}$ & 穿过 $Q$ 的子模型。 \\
$\bar Q_{L_k}$ & $L_k$ 的概率分布的条件均值。 \\
$G_{0,A_k}$ & $A_k$ 的真实条件概率分布。 \\
$G_{A_k}$ & $A_k$ 的可能条件概率分布。 \\
$G_{n,A_k}$ & $G_{0,A_k}$ 的估计。 \\
$\bar Q_n^0$ & $\bar Q_0$ 的初始估计。 \\
$\bar Q_n^1$ & $\bar Q_0$ 的第一次更新估计。 \\
$\bar Q_n^k$ & $\bar Q_0$ 的第 $k$ 次更新估计。 \\
$\bar Q_n^*$ & TMLE 程序中 $\bar Q_0$ 的靶向估计。 \\
$\Psi(Q_0)$ & 当真实目标参数仅通过 $Q_0$ 依赖于 $P_0$ 时，使用的另一种记法。 \\
$\Psi(Q_n^*)$ & 参数的靶向估计量。 \\
$Pf$ & $f(O)$ 在 $P$ 下的期望，例如 $P_0L(Q)=\int L(Q)(o)\,\mathrm dP_0(o)$。 \\
\end{longtable}
